% =====================================================================
% SHARD-ID: LB-03-SYMMETRY-NOETHER
% SHARD-TITLE: Symmetry on a window: intertwiners, bond implementers, and
%              lattice Noether
% SHARD-SOURCES: definitions.md D2, D3, D4, D10 (status header: written
%              2026-08-25, FROZEN at the 2026-08-26 freeze after the Corner A
%              L6 loop converged, theory/verdicts/corner-a-r3.md PASS);
%              notation.md; HANDOFF_MPS_SOFT_THEOREM.md 1.1, 1.2;
%              theory/corner-a.md steps (1.2)-(1.4);
%              theory/corner-a-goldstone.md step (1.6);
%              refs/arxiv-0802.0447, refs/arxiv-2011.12127;
%              theory/checks/corner_a_check.py (checks C0-C2b, C8, C8b)
% SHARD-CLAIMS: D2, D3, D4, D10
% =====================================================================

% --- shard-local semantic macros (mirror notation.md; \providecommand so
% --- that a sibling shard may declare the same symbol without a clash) ---
\providecommand{\Aloc}{\mathfrak{A}_{\mathrm{loc}}}      % local subalgebra
\providecommand{\Utens}{\mathcal{U}}                     % symmetry on a tensor
\providecommand{\Vac}{\Omega_{\mathrm{vac}}}             % vacuum label set
\providecommand{\Nnull}{\mathcal{N}}                     % null-direction map
\providecommand{\Ewin}{\mathfrak{E}}                     % window CP map
\providecommand{\lamt}{\tilde{\lambda}}                  % honest decay rate
\providecommand{\Cbd}{C_{\partial}}                      % window-norm constant
\providecommand{\gaugerem}{\mathfrak{B}}                 % gauge remainder
\providecommand{\Ept}{\mathcal{E}}                       % endpoint space
\providecommand{\dcoset}{\mathfrak{d}}                   % double-coset label
\providecommand{\Fc}{\mathfrak{F}_{c}}                   % finitely supported
\providecommand{\Fec}{\mathfrak{F}_{\mathrm{ec}}}        % eventually constant
\providecommand{\Fell}{\mathfrak{F}_{\ell^{1}}}          % finite variation
\providecommand{\Aeff}{\A_{\mathrm{eff}}}                % effective asg group
\providecommand{\hsub}{\mathfrak{h}}                     % unbroken subalgebra
\providecommand{\msub}{\mathfrak{m}}                     % broken complement
\providecommand{\liealg}{\mathfrak{g}}                   % Lie algebra of G
\providecommand{\Vbond}{\mathcal{V}}                     % bond implementer
\providecommand{\Wspan}{\mathcal{W}}                     % insertion window span
\providecommand{\chalg}{\mathfrak{a}}                    % twisted group algebra
\providecommand{\bcharge}{\mathfrak{q}}                  % asymptotic bond charge
\providecommand{\bondpot}{\mathcal{J}}                   % bond current potential

\section{Symmetry on a window: intertwiners, bond implementers, and lattice
Noether}
\label{sec:symmetry-noether}

An on-site symmetry acts on a matrix-product state in a way that is completely
determined at the level of the single tensor. If the symmetry leaves the state
invariant, then applying the physical unitary $u(g)$ to the physical index of
the tensor cannot change the tensor by more than the residual gauge freedom of
the parametrisation: a phase and a conjugation by some invertible matrix acting
on the virtual (bond) indices. That statement is the \emph{fundamental theorem
of matrix product states}, and the resulting relation --- the
\emph{intertwining relation} of \cref{def:intertwining} --- is the engine that
drives the entire campaign. Because the conjugating matrices $V_{\alpha}(g)$
are determined only up to a phase, they compose only up to a phase: they form
a projective representation, whose cohomology class
$[\omega_{\alpha}] \in H^{2}(H_{\alpha}, U(1))$ is the one-dimensional
symmetry-protected-topological index. On the lattice this class is the exact
analogue of an anomaly of the asymptotic symmetry, and it is where the
campaign's genuinely new prediction lives.

The reason the intertwining relation matters so much is what happens when the
symmetry is applied only to a finite interval $R = [a,b]$ instead of the whole
chain. Substituting the relation at each of the $\abs{R}$ sites of $R$ produces
an alternating string
$V_{\alpha}(g)^{-1} A \, V_{\alpha}(g) V_{\alpha}(g)^{-1} A \, V_{\alpha}(g)
 \cdots$, in which every interior pair
$V_{\alpha}(g) V_{\alpha}(g)^{-1}$ collapses to the identity. Only the two
outermost factors survive, and they sit on the two \emph{boundary bonds} of
$R$. Concretely, with $T_{R}^{(g)}$ denoting the decoration that carries
$A_{\alpha}$ on $\Lambda \setminus R$, $A_{g \cdot \alpha}$ on $R$, the bond
insertion $V_{\alpha}(g)^{-1}$ on $\partial_{-}R = (a-1|a)$ and
$V_{\alpha}(g)$ on $\partial_{+}R = (b|b+1)$,
\begin{equation}
  U_{R}(g) \, \ket{\psi_{\Lambda}(A_{\alpha}; b_{l}, b_{r})}
  \;=\; e^{\,i \abs{R} \theta_{\alpha}(g)} \,
        \ket{\psi_{\Lambda}\bigl(T_{R}^{(g)}; b_{l}, b_{r}\bigr)} ,
  \label{eq:telescoping}
\end{equation}
exactly, as window vectors --- and, as states, with no phase at all, since the
extensive factor $e^{i\abs{R}\theta_{\alpha}(g)}$ appears in bra and ket and
cancels. This is the lattice Ward identity. It says that a symmetry applied
to a region acts \emph{only at the region's boundary}, and acts there by
virtual-space operators that are invisible to any local physical measurement in
the bulk. Every selection rule, every current-conservation statement, and the
whole soft-theorem mechanism of the campaign is this telescoping plus
transfer-matrix calculus. Three points about \eqref{eq:telescoping} deserve
emphasis. The identity is exact on a finite window, provided either that the
window is enlarged to contain $[a-1,b+1]$, or that the window vector admits
insertions on its edge bonds --- which is precisely why
\cref{def:window-vector} was written to admit them. The telescoping uses only
the intertwining relation; injectivity is needed to \emph{obtain} that relation
and to make the pair $(\theta_{\alpha}, V_{\alpha})$ unique, not for the
algebra. And the orientation matters: the surviving operator on the
\emph{left} boundary bond is $V_{\alpha}(g)^{-1}$ and on the \emph{right} it is
$V_{\alpha}(g)$, fixed by the convention of \cref{def:intertwining}. The
original campaign brief states the opposite orientation; that document is
historical and is not edited, so the correction is recorded here. A
$\Z_{2}$ example is blind to the flip, because there $V = Z = V^{-1}$; a
$U(1)$ symmetry with $V(t) = e^{itZ/2}$ separates them cleanly, the correct
orientation giving an error of $5.6 \cdot 10^{-17}$ and the flipped one
$0.267$.

The rest of this section fixes, in order: the symmetry data and the
intertwining relation; the classes of \emph{profiles} $f : \Z \to G$ for which
a modulated symmetry operation makes sense (a discipline that is not
bureaucratic --- half-infinite strings are not elements of the algebra at all);
the bond implementers and the group they generate; and the lattice Noether pair
of charge density and current.

\subsection{The symmetry data and the intertwining relation}

\begin{definition}[On-site symmetry and its action on a tensor (D2, preamble)]
\label{def:onsite-symmetry}
Let $G$ be a compact group (finite, or a compact Lie group) with a unitary
representation $u : G \to U(\C^{d})$ on one site, and let
$U_{\Lambda}(g) := \bigotimes_{x \in \Lambda} u_{x}(g)$. Write $\Utens(g)$ for
the induced linear map on tensors,
\begin{equation}
  \bigl(\Utens(g) A\bigr)^{s} := \sum_{s'} u(g)_{s s'} A^{s'},
  \qquad u(g)_{s s'} := \bra{s} u(g) \ket{s'},
\end{equation}
so that $\Utens(g) \Utens(h) = \Utens(gh)$.
\end{definition}

\begin{definition}[Covariant vacuum family (D2(a))]
\label{def:covariant-vacuum}
A \emph{$G$-covariant vacuum family} consists of a finite or compact set
$\Vac$ of labels, injective canonical-form tensors
$\{A_{\alpha}\}_{\alpha \in \Vac}$ of common bond dimension $\chi$ generating
pairwise distinct states $\omega_{\alpha} := \omega_{A_{\alpha}}$, together
with a $G$-action $(g,\alpha) \mapsto g \cdot \alpha$ on $\Vac$ such that
$U_{\Lambda}(g)$ maps $\omega_{\alpha}$ to $\omega_{g \cdot \alpha}$ for every
finite $\Lambda$, in the sense
\begin{equation}
  \omega_{\alpha} \circ \Ad\bigl(U_{\Lambda}(g)^{\dagger}\bigr)
  \;=\; \omega_{g \cdot \alpha}
  \qquad \text{on } \alg_{\Lambda}, \text{ for all } \Lambda .
\end{equation}
The \emph{stabiliser} $H_{\alpha} := \{ g \in G : g \cdot \alpha = \alpha \}$
is the \emph{unbroken subgroup}. The family is \emph{unbroken at $\alpha$} if
$H_{\alpha} = G$, in which case $\abs{\Vac}$ may be $1$.
\end{definition}

\provenance{D2(a)}%
{\texttt{definitions.md} D2(a); FROZEN at the 2026-08-26 freeze}%
{---}%
{\texttt{theory/verdicts/corner-a-r3.md} (PASS)}

The next relation is the one everything rests on. It is not an assumption of
this campaign: it is the fundamental theorem of matrix product states, and it
is quoted here against the local \TeX{} sources rather than from memory.

\begin{definition}[The intertwining relation (D2(b))]
\label{def:intertwining}
For every $\alpha \in \Vac$ and every $g \in G$ there exist a phase
$\theta_{\alpha}(g) \in \R/2\pi\Z$ and a unitary $V_{\alpha}(g) \in U(\chi)$
with
\begin{equation}
  \sum_{s'} u(g)_{s s'} \, A_{\alpha}^{s'}
  \;=\; e^{\,i \theta_{\alpha}(g)} \,
        V_{\alpha}(g)^{-1} \, A_{g \cdot \alpha}^{s} \, V_{\alpha}(g)
  \qquad \text{for all } s .
  \label{eq:IT}
\end{equation}
The pair $(\theta_{\alpha}(g), V_{\alpha}(g))$ is unique up to
$V_{\alpha}(g) \mapsto e^{i\varphi} V_{\alpha}(g)$.
\end{definition}

\begin{remark}[Ground truth for \eqref{eq:IT}]
Existence: \texttt{refs/arxiv-0802.0447} (\texttt{StringOrder-v10.tex}),
Lemma~1, states that the spectral radius of the mixed transfer map $E_{u}$ is
at most $1$, with equality if and only if there are a unitary $V$ and a phase
$\theta$ with
$V^{\dagger} \tilde{A}_{j} = e^{i(\theta - \theta_{j})} \tilde{A}_{j}
 V^{\dagger}$, together with the accompanying condition
$(u \otimes \mathbb{1}) B = (\mathbb{1} \otimes V) B V^{\dagger}$; and
\texttt{refs/arxiv-2011.12127} (\texttt{TN-Review-main.tex}), which for
$U(g)^{\otimes N} \ket{\psi_{N}} \simeq \ket{\psi_{N}}$ states
$\sum_{j} U_{ij}(g) A^{j} = e^{i\phi(g)} X^{\dagger}(g) A^{i} X(g)$.
Uniqueness: \texttt{refs/arxiv-2011.12127}, the equation labelled
\texttt{XAX=B}, states that
\begin{equation}
  X^{-1} A^{i} X = e^{i\varkappa} \, Y^{-1} A^{i} Y
  \quad \Longrightarrow \quad
  \varkappa = 0 \; \text{ and } \; X = e^{i\varphi} Y
  \text{ for some } \varphi .
\end{equation}
\end{remark}

\provenance{D2(b)}%
{\texttt{definitions.md} D2(b); consequences derived at
 \texttt{theory/corner-a.md} step (1.2)}%
{\texttt{theory/checks/corner\_a\_check.py}, checks C0, C1, C1b, C1c
 (orientation)}%
{\texttt{theory/verdicts/corner-a-r3.md}}

Composing the intertwining relation with itself forces the phase and the
virtual matrices to compose in a definite way. The virtual matrices compose
only projectively, and the class of that projective multiplier is the
symmetry-protected-topological index --- the lattice anomaly.

\begin{definition}[Cocycles and the symmetry-protected index (D2(c))]
\label{def:cocycle-spt-index}
Composition of \eqref{eq:IT} gives, for all $g, h \in G$,
\begin{equation}
  \theta_{\alpha}(hg) = \theta_{\alpha}(g) + \theta_{g \cdot \alpha}(h),
  \qquad
  V_{g \cdot \alpha}(h) \, V_{\alpha}(g)
  = e^{\,i \omega_{\alpha}(h,g)} \, V_{\alpha}(hg) .
  \label{eq:cocycles}
\end{equation}
Restricted to $g,h \in H_{\alpha}$ these say that
$\theta_{\alpha}\vert_{H_{\alpha}} : H_{\alpha} \to U(1)$ is a homomorphism and
that $\omega_{\alpha} : H_{\alpha} \times H_{\alpha} \to U(1)$ is a $2$-cocycle.
Its class
\begin{equation}
  [\omega_{\alpha}] \;\in\; H^{2}\bigl(H_{\alpha}, U(1)\bigr)
\end{equation}
is the \emph{symmetry-protected-topological index} of $A_{\alpha}$.
\end{definition}

\begin{remark}[A notation hazard, stated once]
The symbol $\omega$ carrying \emph{two group} arguments is this cocycle; the
symbol $\omega$ carrying a \emph{momentum} argument, $\omega(k)$, is a magnon
dispersion. The two never occur with the same argument type, and this is the
mechanical disambiguation rule used throughout the repository.
\end{remark}

\begin{definition}[Normal ordering (D2(d))]
\label{def:normal-ordering}
For $g \in H_{\alpha}$ put
$\check{u}_{\alpha}(g) := e^{-i\theta_{\alpha}(g)} u(g)$. Because
$\theta_{\alpha}\vert_{H_{\alpha}}$ is a homomorphism, $\check{u}_{\alpha}$ is
again a unitary representation of $H_{\alpha}$, and it satisfies
\eqref{eq:IT} with phase $1$. All charges below are those of
$\check{u}_{\alpha}$ unless stated otherwise; physically this is the charge
measured from its vacuum value.
\end{definition}

\begin{definition}[Smoothness hypothesis (S) (D2(e))]
\label{def:smoothness}
When $G$ is a Lie group we assume, and state whenever it is used, that
$\varepsilon \mapsto A_{\exp(\varepsilon\xi) \cdot \alpha}$ and
$\varepsilon \mapsto V_{\alpha}(\exp(\varepsilon\xi))$ admit $C^{1}$ local
sections with $V_{\alpha}(e) = \mathbb{1}$. We then write
\begin{equation}
  X_{\alpha}(\xi) := \frac{d}{d\varepsilon}
     V_{\alpha}\bigl(\exp(\varepsilon\xi)\bigr)\Big\vert_{\varepsilon = 0}
  \quad \text{(anti-Hermitian)},
  \qquad
  \theta'_{\alpha}(\xi) := \frac{d}{d\varepsilon}
     \theta_{\alpha}\bigl(\exp(\varepsilon\xi)\bigr)\Big\vert_{\varepsilon = 0}.
\end{equation}
\end{definition}

\provenance{D2(c),(d),(e)}%
{\texttt{definitions.md} D2(c)--(e); the cocycle relations
 \eqref{eq:cocycles} derived at \texttt{theory/corner-a.md} step (1.2)}%
{\texttt{theory/checks/corner\_a\_check.py}}%
{\texttt{theory/verdicts/corner-a-r3.md}}

\subsection{Which profiles are allowed}

A ``modulated symmetry'' means applying $u(f(x))$ at site $x$ for some profile
$f : \Z \to G$. Whether the resulting object is an operator, a map on states,
or nothing at all depends entirely on how $f$ behaves at infinity. The
following classes are not bookkeeping: a half-infinite symmetry string is
\emph{not} an element of the quasi-local algebra, and a plane wave is in none
of the classes and is admissible only inside a Fourier transform. Every soft
statement in the campaign must say which class it uses.

\begin{definition}[Admissible profile classes (D3(a))]
\label{def:admissible-profiles}
Let
\begin{align}
  \Fc  &:= \bigl\{\, f : \Z \to G \;:\; f(x) = e
           \text{ for all but finitely many } x \,\bigr\}, \\
  \Fec &:= \bigl\{\, f : \Z \to G \;:\; f \text{ is constant on some }
           (-\infty,-n] \text{ and on some } [n,\infty) \,\bigr\},
\end{align}
and, for a Lie algebra element $\xi \in \liealg$,
\begin{align}
  \Fc(\xi)   &:= \bigl\{\, f : \Z \to \R \;:\; f
                 \text{ finitely supported} \,\bigr\}, \\
  \Fec(\xi)  &:= \bigl\{\, f : \Z \to \R \;:\; f
                 \text{ constant near } \pm\infty \,\bigr\}, \\
  \Fell(\xi) &:= \Bigl\{\, f : \Z \to \R \;:\;
                 \sum_{x} \abs{f(x+1) - f(x)} < \infty \,\Bigr\}
                 \quad \text{(finite total variation)} .
\end{align}
One has $\Fc \subset \Fec$ and $\Fec(\xi) \subset \Fell(\xi)$. A profile is
\emph{admissible for an operator statement} if and only if it lies in $\Fc$,
in which case the operator lies in $\Aloc$; it is \emph{admissible for a state
statement} if and only if it lies in $\Fec$, in which case the operator is a
weak-$*$ limit of elements of $\Aloc$. The plane wave $f(x) = e^{ikx}$ lies in
none of these for $k \neq 0$; it is admitted only inside a Fourier transform,
i.e.\ as the distributional kernel of a wave packet
$\sum_{k} \varphi(k) f_{k}$ with $\varphi \in C_{c}^{\infty}((-\pi,\pi])$.
\end{definition}

\begin{definition}[Truncated symmetry operations (D3(b))]
\label{def:truncated-symmetry}
For a finite interval $R = [a,b] \subset \Z$ put
$U_{R}(g) := \prod_{x \in R} u_{x}(g) \in \Aloc$. Its \emph{boundary bonds}
are $\partial_{-}R := (a-1|a)$ and $\partial_{+}R := (b|b+1)$, and
$\abs{R} = b-a+1$. More generally, for $f \in \Fc$,
$U[f] := \prod_{x} u_{x}(f(x)) \in \Aloc$.
\end{definition}

\begin{definition}[Half-infinite strings act on states only (D3(c))]
\label{def:half-infinite-string}
For $f \in \Fec$ the product $U[f]$ is \emph{not} an element of $\alg$. It is
used only through the induced map on states,
\begin{equation}
  \varrho \;\longmapsto\;
  \text{w*-}\!\lim_{n \to \infty} \varrho \circ \Ad\bigl(U[f_{n}]^{\dagger}\bigr),
\end{equation}
where $f_{n} \in \Fc$ agrees with $f$ on $[-n,n]$ and equals $e$ outside.
Existence of this limit is a theorem, not a definition; when it exists the
limit state is written $f \triangleright \varrho$. The canonical case is
$f = 1_{[x,\infty)} \cdot g$, written $U_{[x,\infty)}(g)$.
\end{definition}

\provenance{D3}%
{\texttt{definitions.md} D3(a)--(c); existence of the half-infinite limit
 proved at \texttt{theory/corner-a.md} steps (1.2)(c) and (1.4)}%
{---}%
{\texttt{theory/verdicts/corner-a-r3.md}}

\subsection{Bond implementers and the effective asymptotic symmetry group}

The telescoping identity \eqref{eq:telescoping} leaves virtual operators on
bonds. To speak of them as \emph{operators} one needs the map
``$M \mapsto$ the window vector with $M$ inserted on bond $b_{0}$'' to be
injective, so that left multiplication on $M$ transports to a well-defined map
on window vectors. That injectivity is genuinely conditional: it needs enough
uniform sites on each side of the bond, and without them the construction is
ill defined --- not merely unproved. The definition therefore carries its
padding hypothesis explicitly.

\begin{definition}[Padded windows and the insertion map (D4(a1),(a2))]
\label{def:padded-window}
A window $\Lambda = [a,b]$ is \emph{padded about the bond $b_{0} = (m|m+1)$}
if it contains at least $n_{0}$ sites on each side of $b_{0}$, with $n_{0}$ the
word-spanning length of \eqref{eq:word-spanning}, and if the boundary vectors
satisfy $b_{l} \neq 0$ and $b_{r} \neq 0$. For such $\Lambda$ let
$\Wspan_{\Lambda, b_{0}}$ be the linear span of the window vectors carrying one
insertion on $b_{0}$, and let
\begin{equation}
  \iota_{\Lambda, b_{0}} : M_{\chi}(\C) \longrightarrow
  \Wspan_{\Lambda, b_{0}},
  \qquad
  M \longmapsto \ket{\psi_{\Lambda}(A; b_{l}, b_{r}; M@b_{0})} .
\end{equation}
On a padded window, $\iota_{\Lambda,b_{0}}$ is \textbf{injective}.
\end{definition}

\begin{proof}[Idea of proof of injectivity]
The coefficient of $\ket{s}$ is $b_{l}^{\dagger} P(s) \, M \, Q(s) \, b_{r}$,
with $P$ and $Q$ the matrix words to the left and to the right of $b_{0}$. By
\eqref{eq:word-spanning} each of $P$, $Q$ ranges over a spanning set of
$M_{\chi}(\C)$, and the coefficient is linear in each; since $b_{l} \neq 0$
the set $\{ b_{l}^{\dagger} P : P \in M_{\chi} \}$ is all of
$(\C^{\chi})^{*}$, and since $b_{r} \neq 0$ the set $\{ Q b_{r} \}$ is all of
$\C^{\chi}$. So $\iota(M) = \iota(M')$ forces $v^{\dagger}(M-M')w = 0$ for all
$v, w$, i.e.\ $M = M'$.
\end{proof}

\begin{remark}[Padding is necessary, not decorative]
\label{rem:padding-necessary}
Without padding the kernel of $\iota$ need not be invariant under left
multiplication, so the rule ``$\iota(M) \mapsto \iota(NM)$'' is ill defined.
An explicit injective, canonical-gaugeable, $\Z_{2}$-symmetric counterexample:
take
\begin{equation}
  A^{0} = \operatorname{diag}(1,2), \quad A^{1} = X, \quad
  u(g) = \operatorname{diag}(1,-1), \quad V = Z, \quad
  b_{l} = (\sqrt{2},\, 1), \quad b_{r} = (1,\, 0),
\end{equation}
one site on each side of the bond, and
\begin{equation}
  N = \begin{pmatrix} -\sqrt{2} & 0 \\ 1 & 0 \end{pmatrix} .
\end{equation}
Then $\iota(N) = 0$ while $\norm{\iota(ZN)}_{\infty} = 4$. The transfer
spectrum is $\{4.303,\, 3,\, 1,\, 0.697\}$ --- a unique top eigenvalue, with
$\lambda_{E} = 0.697$ after rescaling --- and the length-two words have rank
$4$, so this genuinely is a tensor of the class of
\cref{def:injective-mps}. Padding both sides to $n_{0} = 2$ restores
$\ran \iota$ of full dimension $\chi^{2} = 4$.
\end{remark}

\begin{definition}[Bond implementers (D4(a4))]
\label{def:bond-implementer}
On a padded window define, for $M \in GL(\chi)$,
\begin{equation}
  \Vbond_{b_{0}}(M) \; := \;
  \iota_{\Lambda,b_{0}} \circ L_{M} \circ \iota_{\Lambda,b_{0}}^{-1}
  \qquad \text{on } \ran \iota_{\Lambda,b_{0}} = \Wspan_{\Lambda,b_{0}},
\end{equation}
where $L_{M}$ is left multiplication on $M_{\chi}(\C)$. By
\cref{def:padded-window} this is well defined and linear, and
$\Vbond_{b_{0}}(M)\Vbond_{b_{0}}(M') = \Vbond_{b_{0}}(MM')$, with scalars
acting by scalars. Put
\begin{equation}
  \Vbond_{b_{0}}(g) \; := \; \Vbond_{b_{0}}\bigl(V_{\alpha}(g)\bigr)
  \qquad \text{(\textbf{not} } V_{\alpha}(g)^{-1} \text{)} .
\end{equation}
With the convention \eqref{eq:cocycles}, this gives
\begin{equation}
  \Vbond_{b_{0}}(h) \, \Vbond_{b_{0}}(g)
  \;=\; e^{\,i \omega_{\alpha}(h,g)} \, \Vbond_{b_{0}}(hg) ,
  \label{eq:implementer-cocycle}
\end{equation}
a linear representation of the twisted group algebra $\chalg_{\alpha}$ of
\cref{def:charge-algebra} on $\Wspan_{\Lambda,b_{0}}$, in which the multiplier
acts nontrivially. On \emph{states} the multiplier is invisible, since
$\omega^{cM@b} = \omega^{M@b}$ for a scalar $c$; the state-level action is
\cref{def:effective-asg}.
\end{definition}

\begin{remark}[Two corrections carried in this definition]
The r1 version used the implementer $V_{\alpha}(g)^{-1}$ with left
multiplication, which reverses the composition law, and asserted a
twisted-algebra action on states, where phases are invisible. The r2 version
declared the implementer on an unpadded insertion span, where it is not well
defined --- see \cref{rem:padding-necessary}. Both are superseded by the form
above. This definition applies in the unbroken case $H_{\alpha} = G$; the
broken case is handled through \cref{def:kink-sector} and the Corner~A
endpoint theorem.
\end{remark}

\provenance{D4(a)}%
{\texttt{definitions.md} D4(a1)--(a4); consumed at \texttt{theory/corner-a.md}
 step (1.4)}%
{\texttt{theory/checks/corner\_a\_check.py}, checks C8 (kernel not invariant
 without padding) and C8b (padding to $n_{0}=2$ restores injectivity)}%
{\texttt{theory/verdicts/corner-a-r2.md} objections 1, 3, 4, 7, 8;
 \texttt{theory/verdicts/corner-a-r3.md}}

Passing from window vectors to states loses exactly the phases, and with them
the multiplier. What survives is an honest group action, and the group that
acts is not $G$ but $G$ modulo those elements whose virtual matrix is a scalar.

\begin{definition}[The effective asymptotic symmetry group (D4(b))]
\label{def:effective-asg}
Set
\begin{equation}
  N_{\alpha} \; := \; \bigl\{\, g \in G \;:\;
  V_{\alpha}(g) \in \C^{\times}\mathbb{1} \,\bigr\},
\end{equation}
a normal subgroup of $G$ --- the preimage of the centre under a projective
representation. The induced map on states is the group homomorphism
\begin{equation}
  \rho_{\alpha} : G \longrightarrow PGL(\chi),
  \qquad \rho_{\alpha}(g) := [V_{\alpha}(g)],
\end{equation}
which is \emph{not} projective, the cocycle having been quotiented away, and
which has $\ker \rho_{\alpha} = N_{\alpha}$. The \emph{effective asymptotic
symmetry group} is
\begin{equation}
  \Aeff \; := \; G/N_{\alpha} \;\cong\; \rho_{\alpha}(G) \;\subseteq\; PGL(\chi) .
  \label{eq:aeff}
\end{equation}
This is a genuine group.
\end{definition}

\begin{remark}[The superseded object]
\label{rem:asg-superseded}
The object $\A = (G_{L} \times G_{R})/G_{\mathrm{diag}}$ of the original brief
is \emph{not} the vacuum orbit unless $N_{\alpha} = \{e\}$, and it is a group
only for abelian $G$. It is retained only as a name for the coset space and is
never used as a classifying object; the classifying objects are $\Aeff$ in the
unbroken case and the double coset \eqref{eq:double-coset} in the broken case.
\end{remark}

\begin{definition}[The two asymptotic copies (D4(c))]
\label{def:two-asymptotic-copies}
$G_{L}$ and $G_{R}$ act through half-infinite strings on $(-\infty,b]$ and
$[b+1,\infty)$ respectively. By the telescoping identity
\eqref{eq:telescoping}, the left string leaves $V_{\alpha}(g_{L})$ on the bond
$b$ and the right string leaves $V_{\alpha}(g_{R})^{-1}$, so the composite
residue is $V_{\alpha}(g_{L}) V_{\alpha}(g_{R})^{-1}$, which is trivial on
states if and only if $g_{L} g_{R}^{-1} \in N_{\alpha}$. Hence the true
stabiliser is
\begin{equation}
  S_{\alpha} := \bigl\{ (g_{L},g_{R}) \;:\;
  g_{L} g_{R}^{-1} \in N_{\alpha} \bigr\} \;\supseteq\; G_{\mathrm{diag}},
\end{equation}
with $S_{\alpha} = G_{\mathrm{diag}}$ if and only if $N_{\alpha} = \{e\}$.
\end{definition}

\begin{definition}[Asymptotic charge algebra (D4(d))]
\label{def:charge-algebra}
$\chalg_{\alpha} := \C_{\omega_{\alpha}}[H_{\alpha}]$ is the twisted group
algebra with basis $\{v_{g}\}$ and product
$v_{h} v_{g} = e^{i\omega_{\alpha}(h,g)} v_{hg}$. This definition supplies
only the name; the content --- that $g \mapsto V_{\alpha}(g)$ is a
$*$-representation of it, and that $\rho_{\alpha}$ is the induced state-level
action --- is a theorem of the Corner~A endpoint shard. For a Lie group,
under the smoothness hypothesis of \cref{def:smoothness}, the
\emph{asymptotic charges} are $\bcharge_{b}(\xi) := X_{\alpha}(\xi)$, acting by
left multiplication on the bond insertion, and
\begin{equation}
  \bigl[\, \bcharge_{b}(\xi), \, \bcharge_{b}(\zeta) \,\bigr]
  \;=\; \bcharge_{b}\bigl([\xi,\zeta]\bigr)
        \; + \; c_{\alpha}(\xi,\zeta) \, \mathbb{1} .
  \label{eq:charge-algebra}
\end{equation}
\end{definition}

\begin{remark}[What the central term is, and what it is not]
\label{rem:no-central-charge}
The term $c_{\alpha}$ in \eqref{eq:charge-algebra} is only the \emph{local
infinitesimal image} of the class $[\omega_{\alpha}]$, and it generally loses
it. For $\hsub_{\alpha}$ compact semisimple, Whitehead's second lemma says
that its Lie-algebra cohomology class in $H^{2}(\hsub_{\alpha},\R)$ is
trivial; the displayed central cocycle is then a coboundary and can be gauged
away by choosing a phase section, although it need not vanish in every local
section. The global class $[\omega_{\alpha}]$ may nevertheless be nontrivial
torsion: for the AKLT chain, $H^{2}(SO(3),U(1)) = \Z_{2}$, while a suitable
phase convention gives the ordinary $su(2)$ bracket for the spin-one-half edge
generators. \textbf{The lattice symmetry-protected anomaly is a
group-cohomological multiplier, not a Lie-algebra central charge.} This is a
genuine disanalogy with the continuum picture of a centrally extended charge
algebra, and it must be stated as such rather than glossed over.
\end{remark}

\provenance{D4(b),(c),(d)}%
{\texttt{definitions.md} D4(b)--(d); the representation content proved at
 \texttt{theory/corner-a.md} step (1.4), sub-steps (d) and (f)}%
{\texttt{theory/checks/corner\_a\_check.py}}%
{\texttt{theory/verdicts/corner-a-r2.md} objections 3, 4;
 \texttt{theory/verdicts/corpus-r2.md} note N5 (the central-term caveat);
 \texttt{theory/verdicts/corner-a-r3.md}}

\subsection{The lattice Noether pair}

Differentiating the on-site symmetry gives a charge density; commuting the
Hamiltonian with the half-line charge gives a current supported near the cut.
The pair satisfies an exact lattice continuity equation --- no continuum limit,
no approximation. The last clause is the one that connects back to the
telescoping mechanism: acting on the vacuum, the physical charge density is the
lattice divergence of a purely virtual, bond-supported quantity. That
quantity is the lattice analogue of a gauge potential, and the charge density
is its field strength.

\begin{definition}[Charge density and cut current (D10(a),(b))]
\label{def:noether-pair}
Write $\liealg = \mathrm{Lie}(G)$, and let $H = \sum_{x} h_{x,x+1}$ be a
translation-invariant nearest-neighbour Hamiltonian; the definition extends
verbatim to finite range $R_{h}$. For $\xi \in \liealg$ let
\begin{equation}
  q(\xi) := \frac{d}{d\varepsilon} u\bigl(\exp(\varepsilon\xi)\bigr)
  \Big\vert_{\varepsilon=0}
  \quad \text{(anti-Hermitian on } \C^{d} \text{)},
\end{equation}
with $q_{x}(\xi)$ its copy at site $x$. Assume $H$ is $G$-invariant on-site:
$[\,h_{x,x+1},\, q_{x}(\xi) + q_{x+1}(\xi)\,] = 0$.

\emph{(a) Cut current.} For finite range $R_{h}$, each $h_{x}$ supported in
$[x, x+R_{h}-1]$, and $G$-invariance
$[\,h_{x}, \sum_{y \in \supp h_{x}} q_{y}(\xi)\,] = 0$, define the current
across the bond $(m|m+1)$ as
\begin{equation}
  j_{m|m+1}(\xi) \; := \; -\Bigl[\, H, \sum_{y \le m} q_{y}(\xi) \,\Bigr]
  \; = \; -\sum_{x} \Bigl[\, h_{x}, \sum_{y \le m} q_{y}(\xi) \,\Bigr].
  \label{eq:cut-current}
\end{equation}
The sum over $x$ is finite --- at most $R_{h}$ terms, those whose support
straddles the cut --- because $[\,h_{x}, \sum_{y \le m} q_{y}\,] = 0$ whenever
$\supp h_{x}$ lies wholly on one side. Hence
$j_{m|m+1}(\xi) \in \Aloc$. In the nearest-neighbour case $R_{h} = 2$ this
reduces to
\begin{equation}
  j_{x|x+1}(\xi) = -\bigl[\, h_{x,x+1}, q_{x}(\xi) \,\bigr]
                 = \bigl[\, h_{x,x+1}, q_{x+1}(\xi) \,\bigr].
\end{equation}

\emph{(b) Continuity equation.} Exactly, on the lattice,
\begin{equation}
  \bigl[\, H, q_{x}(\xi) \,\bigr]
  \;=\; j_{x-1|x}(\xi) \; - \; j_{x|x+1}(\xi) .
  \label{eq:continuity}
\end{equation}
\end{definition}

\begin{definition}[Modulated charge and current (D10(c))]
\label{def:modulated-charge}
For $f \in \Fc(\xi)$,
\begin{equation}
  Q[f;\xi] := \sum_{x} f(x) \, q_{x}(\xi) \in \Aloc,
  \qquad
  J[f;\xi] := \sum_{x} f(x) \, j_{x|x+1}(\xi) .
\end{equation}
For $f(x) = e^{ikx}$ we write $Q_{k}(\xi)$ and $J_{k}(\xi)$, understood in the
wave-packet sense of \cref{def:admissible-profiles}.
\end{definition}

\begin{definition}[Bond current potential (D10(d))]
\label{def:bond-current-potential}
In the unbroken case, with the intertwining relation \eqref{eq:IT} and normal
ordering (\cref{def:normal-ordering}), let $\bondpot_{b}(\xi)$ denote the
operation ``insert $X_{\alpha}(\xi)$ on bond $b$'', in the sense of
\cref{def:bond-implementer}. The Goldstone shard's theorem G0(d) states that
\begin{equation}
  q_{x}(\xi) \triangleright \omega_{\alpha}
  \;=\; \bigl( \bondpot_{x|x+1}(\xi) - \bondpot_{x-1|x}(\xi) \bigr)
        \triangleright \omega_{\alpha} :
  \label{eq:bond-potential}
\end{equation}
the physical charge density acting on the vacuum is the lattice divergence of a
purely virtual, bond-supported quantity. In this dictionary $\bondpot$ is the
lattice analogue of the gauge potential and $q$ is the field strength.
\end{definition}

\begin{remark}[A definition that was repaired]
The r1 text defined the current by the nearest-neighbour formula while the
hypothesis it served quantified over finite range. The cut-current form
\eqref{eq:cut-current} is the repaired definition, and it supersedes the
nearest-neighbour one, which remains correct in its original restricted
setting.
\end{remark}

\provenance{D10}%
{\texttt{definitions.md} D10(a)--(d); the continuity equation
 \eqref{eq:continuity} and the bond-potential identity
 \eqref{eq:bond-potential} proved at
 \texttt{theory/corner-a-goldstone.md} step (1.6)}%
{\texttt{theory/checks/corner\_a\_check.py}}%
{\texttt{theory/verdicts/corner-a-r2.md} objection 19;
 \texttt{theory/verdicts/corner-a-r3.md}}
